
\chapter*{Preámbulo}
\thispagestyle{empty}
    La industria de la construcción de edificios consiste en un proceso que requiere la comunicación de diferentes agentes involucrados en un proyecto, desde el momento en el que se diseñan los planos hasta que el edificio está listo. Arquitectos, ingenieros, contratistas y clientes deben compartir una visión común del edificio que, tradicionalmente, se ha transmitido mediante planos en papel o bidimensionales, junto con maquetas o imágenes estáticas que no siempre facilitan una compresión espacial completa del proyecto.

    Esta desconexión entre la representación técnica y la percepción real genera problemas recurrentes en esta industria: malinterpretaciones de diseño, errores no detectados hasta fases avanzadas de obra, dificultades para validar decisiones en la construcción y barreras en la comunicación con clientes que no tienen conocimientos técnicos. Las metodologías actuales, aunque rigurosas, no explotan el potencial de las tecnologías de realidad inmersiva, las cuales podrían transformar cómo se visualiza y verifica un proyecto arquitectónico.

    Este trabajo propone una solución con el desarrollo de una aplicación que integre Realidad Virtual y/o Realidad Aumentada en el flujo de trabajo de construcción, permitiendo visualizar modelos \gls{bim} directamente en éstas. La investigación previa identificará las necesidades específicas del sector y evaluará qué tecnologías y enfoques resultan más efectivos para cada tipo de tarea.

    El resultado será una aplicación multiplataforma que combine visualización 3D en tiempo real integración con modelos BIM existentes y funcionalidades adaptadas a casos de uso concretos como inspecciones de obra, formación de seguridad o presentaciones a clientes, respaldada por una arquitectura técnica robusta capaz de gestionar grandes volúmenes de información geométrica.


\chapter*{Motivación, justificación y objetivo general}

    Hace alrededor de 2 años, compré mis primeras gafas de realidad virtual, las Meta Quest 2, y es algo de lo que no me arrepiento para nada haber obtenido, ya que es una experiencia que opino que todo el mundo debería experimentar al menos una vez, la inmersión que te da pareciendo que estás en otro mundo da otra sensación de control, sobre todo en los videojuegos.
    Mientras reflexionaba sobre qué tema escogería para mi \gls{tfg}, llegué a la conclusión de que la industria de la realidad virtual o aumentada tiene mucho futuro y margen de mejora, las gafas que tengo, aunque no sean las más recientes, siguen siendo modernas, pero tienen problemas que necesitan arreglo, como el sobrecalentamiento, la potencia o su pesado diseño. Problemas que el siguiente modelo soluciona ligeramente pero no al completo.
    Por razones como esta pienso que los campo de la \gls{ar} y \gls{vr} son de las más interesantes para estudiar ya que, al igual que la \gls{ia}, en un futuro tendrán mucha más relevancia en el ámbito lúdico y laboral, siendo en este caso la construcción de edificios, una de las más imprescindibles en la sociedad. En un tiempo se perfeccionarán estos dispositivos de realidad virtual con el objetivo de que sean más accesibles, versátiles y normalizados.
    Aunque en la formación obtenida actualmente esta industria no es algo en lo que haya llegado a profundizar mucho, me ha ayudado a asentar las bases de la programación, gestión de proyectos y diseño de interfaces que me ayudarán para realizar el proyecto. Este TFG lo utilizaré como entrada para especializarme en este campo.

\cleardoublepage %salta a nueva página impar
\chapter*{Agradecimientos}

\thispagestyle{empty}
\vspace{1cm}

    Con este trabajo me gustaría agradecer a los familiares, empezando por mi madre y mi padrastro por el apoyo incondicional que me han dado por todas mis difíciles etapas de estudio, ayudándome a salir adelante y levantándome cada vez que lo necesitaba para seguir esforzándome.
    Además también quiero dar las gracias a mi hermano, que me ha echado una mano con las asignaturas que he tenido dificultad durante la carrera y me ha enseñado que  llegar hasta aquí siempre ha sido posible.
    Por otro lado también quisiera mostrar agradecimiento a mis amigos, sobre todo a los que se han quedado conmigo hasta el final de esta dura carrera, por escucharme en cada momento y crear buenos momentos, los cuales hasta el día de hoy siguen siendo dignos de recordar.
    Por último, quiero agradecer también a mis compañeros del grupo de \gls{abp}, por haber hecho de ese proyecto incluso de esa magnitud un proceso ameno y en el que me lo he llegado a pasar bien.

    

\cleardoublepage %salta a nueva página impar

\begin{quote}
``\textit{Listen, it's the times when you're scared or worried that you should deal with smiling!}''

\raggedleft --- All Might (Kenta Miyake)
\end{quote}

\begin{quote}
``\textit{I will fight for my future! I will fight alongside the future that everyone created for me! It’s not for anyone else’s sake but my own!}''

\raggedleft --- Hajime Hinata (Johnny Yong)
\end{quote}


%\cleardoublepage %salta a nueva página impar
% Aquí va la cita célebre si la hubiese. Si no, comentar la(s) linea(s) siguientes
%\chapter*{}
%\setlength{\leftmargin}{0.5\textwidth}
%\setlength{\parsep}{0cm}
%\addtolength{\topsep}{0.5cm}
%\begin{flushright}
%\small\em{
%Si consigo ver más lejos\\
%es porque he conseguido auparme\\ 
%a hombros de gigantes
%}
%\end{flushright}
%\begin{flushright}
%\small{
%Isaac Newton.
%}
%\end{flushright}
%\cleardoublepage %salta a nueva página impar
